\documentclass[12pt,a4paper]{article}

\usepackage[utf8]{inputenc}
\usepackage[T1]{fontenc}
\usepackage[margin=2.5cm]{geometry}
\usepackage{booktabs}
\usepackage{longtable}
\usepackage{array}
\usepackage{listings}
\usepackage{xcolor}
\usepackage{hyperref}
\usepackage{parskip}
\usepackage{caption}
\usepackage{titlesec}
\usepackage{enumitem}

\hypersetup{
    colorlinks=true,
    linkcolor=black,
    urlcolor=blue
}

\lstdefinestyle{python}{
    language=Python,
    basicstyle=\ttfamily\small,
    keywordstyle=\color{blue}\bfseries,
    commentstyle=\color{gray},
    stringstyle=\color{orange!70!black},
    numberstyle=\tiny\color{gray},
    numbers=left,
    stepnumber=1,
    breaklines=true,
    frame=single,
    backgroundcolor=\color{gray!5},
    showstringspaces=false,
    tabsize=4,
}

\lstdefinestyle{sql}{
    language=SQL,
    basicstyle=\ttfamily\small,
    keywordstyle=\color{blue}\bfseries,
    commentstyle=\color{gray},
    stringstyle=\color{orange!70!black},
    breaklines=true,
    frame=single,
    backgroundcolor=\color{gray!5},
    showstringspaces=false,
    morekeywords={COPY,STDIN,BIGSERIAL,SMALLINT,SERIAL,BOOLEAN,DECIMAL,VARCHAR,CHAR,REFERENCES,CASCADE,ON,CONFLICT,DO,NOTHING},
}

\title{\textbf{Laboratory 5 Report}\\[0.5em]
\large ETL Process for the Flight Delays Data Warehouse}
\date{June 2026}

\begin{document}

\maketitle
\tableofcontents
\newpage

% ===========================================================================
\section{Analysis of Data Sources}
% ===========================================================================

Three datasets were considered for this project. Before beginning the ETL
design, each was analysed for quality issues, structural incompatibilities,
and missing data.

% ---------------------------------------------------------------------------
\subsection{Dataset 3 — Primary Flight Data (2015)}
% ---------------------------------------------------------------------------

Dataset~3 is the primary source and consists of three CSV files located in
\texttt{data/dataset3/}:

\begin{itemize}
    \item \textbf{\texttt{flights.csv}} — approximately 5.8~million rows of US
          domestic flight records for the year 2015, with 31 columns covering
          scheduling, delays, cancellations, and operational metrics.
    \item \textbf{\texttt{airlines.csv}} — 14~rows mapping IATA carrier codes
          to full airline names.
    \item \textbf{\texttt{airports.csv}} — 322~rows with IATA airport codes,
          names, cities, states, and geographic coordinates.
\end{itemize}

The following quality issues were identified and resolved:

\begin{longtable}{>{\raggedright\arraybackslash}p{4.5cm}
                  >{\raggedright\arraybackslash}p{3.5cm}
                  >{\raggedright\arraybackslash}p{2cm}
                  >{\raggedright\arraybackslash}p{4.5cm}}
\toprule
\textbf{Issue} & \textbf{Column(s)} & \textbf{Rows affected} & \textbf{Resolution} \\
\midrule
\endfirsthead
\toprule
\textbf{Issue} & \textbf{Column(s)} & \textbf{Rows affected} & \textbf{Resolution} \\
\midrule
\endhead
\bottomrule
\endfoot

Missing departure and arrival delays &
\texttt{DEPARTURE\_DELAY}, \texttt{ARRIVAL\_DELAY} &
{\raise0pt\hbox{$\sim$}}86\,153 &
Expected for cancelled flights; stored as \texttt{NULL} \\[4pt]

Missing tail number &
\texttt{TAIL\_NUMBER} &
{\raise0pt\hbox{$\sim$}}14\,721 &
Mapped to \texttt{UNKNOWN} sentinel row in \texttt{Dim\_Plane} \\[4pt]

All delay breakdown columns empty &
\texttt{AIR\_SYSTEM\_DELAY} through \texttt{WEATHER\_DELAY} &
{\raise0pt\hbox{$\sim$}}4\,755\,640 &
Stored as \texttt{0} — absence of delay is not a defect \\[4pt]

Cancellation code absent on non-cancelled flights &
\texttt{CANCELLATION\_REASON} &
{\raise0pt\hbox{$\sim$}}5\,729\,195 &
Mapped to sentinel code \texttt{N} (Not Applicable) in \texttt{Dim\_Cancellation\_Reason} \\[4pt]

Times stored as HHMM integer strings &
\texttt{SCHEDULED\_DEPARTURE}, \texttt{DEPARTURE\_TIME}, \texttt{ARRIVAL\_TIME}, \texttt{SCHEDULED\_ARRIVAL} &
All rows &
Converted to \texttt{HH:MM:00} format; edge case \texttt{2400} mapped to \texttt{00:00:00} \\[4pt]

Date split across three separate columns &
\texttt{YEAR}, \texttt{MONTH}, \texttt{DAY} &
All rows &
Reconstructed into a single \texttt{DATE} value and integer \texttt{date\_key} \\[4pt]

No aircraft metadata &
\texttt{TAIL\_NUMBER} only &
All rows &
\texttt{Dim\_Plane} stores tail numbers only; \texttt{model}, \texttt{manufacturer}, and \texttt{issue\_date} remain \texttt{NULL} \\

\end{longtable}

A specific technical issue was encountered with HHMM time columns: because
values such as \texttt{0005} or \texttt{0930} could be parsed by pandas as
integers, leading zeroes would be silently dropped. This was prevented by
reading all time columns with explicit \texttt{dtype=str} in the extraction
layer.

% ---------------------------------------------------------------------------
\subsection{Dataset 1 — Supplementary Flight Data (2019–2023)}
% ---------------------------------------------------------------------------

Dataset~1 (\texttt{data/dataset1/flights\_sample\_3m.csv}) contains
approximately 3~million rows covering the years 2019 through~2023. Although
it records the same domain (US domestic flights), its schema differs from
Dataset~3, requiring normalisation before it can be fed into the same
transformation pipeline.

\begin{longtable}{>{\raggedright\arraybackslash}p{4.5cm}
                  >{\raggedright\arraybackslash}p{3.5cm}
                  >{\raggedright\arraybackslash}p{2cm}
                  >{\raggedright\arraybackslash}p{4.5cm}}
\toprule
\textbf{Issue} & \textbf{Column(s)} & \textbf{Rows affected} & \textbf{Resolution} \\
\midrule
\endfirsthead
\toprule
\textbf{Issue} & \textbf{Column(s)} & \textbf{Rows affected} & \textbf{Resolution} \\
\midrule
\endhead
\bottomrule
\endfoot

Different column naming conventions &
All columns &
All rows &
Renamed to unified internal names in \texttt{normalize\_dataset1\_chunk()} \\[4pt]

No \texttt{TAIL\_NUMBER} column &
— &
All rows &
\texttt{plane\_key} set to the \texttt{UNKNOWN} sentinel for all Dataset~1 rows \\[4pt]

Airlines not in \texttt{airlines.csv} &
\texttt{AIRLINE\_CODE} &
Varies &
Pre-scan of Dataset~1 inserts any missing carriers into \texttt{Dim\_Airline} before fact loading \\[4pt]

Flight date as \texttt{YYYY-MM-DD} string &
\texttt{FL\_DATE} &
All rows &
Parsed directly; no HHMM conversion required \\[4pt]

Delay breakdown columns differently named &
e.g.\ \texttt{DELAY\_DUE\_CARRIER} vs \texttt{AIRLINE\_DELAY} &
All rows &
Unified by column rename during normalisation \\

\end{longtable}

% ---------------------------------------------------------------------------
\subsection{Dataset 2 — Excluded}
% ---------------------------------------------------------------------------

Dataset~2 (\texttt{data/dataset2/airline.csv.shuffle}, 11~GB) was evaluated
and excluded from the ETL for the following reasons:

\begin{enumerate}
    \item \textbf{File size.} 11~GB exceeds practical processing limits for a
    local academic environment without a distributed computing framework.
    \item \textbf{Encoding.} The file uses latin-1 (non-UTF-8) encoding, which
    creates normalisation complexity and risks mojibake in string columns.
    \item \textbf{Missing reference data.} No airport or airline lookup tables
    accompany the file, making it impossible to enrich records with the
    geographic or carrier metadata required by the dimensional model.
    \item \textbf{No data dictionary.} Some column semantics are ambiguous
    without formal documentation.
    \item \textbf{Sufficient coverage from other sources.} Datasets~1 and~3
    together cover approximately 8.8~million flight records spanning 2015
    and 2019–2023, which provides adequate volume and temporal range for all
    planned analyses.
\end{enumerate}

The only artefact retained from Dataset~2 is
\texttt{data/dataset2/carriers.csv} (1\,492 airline codes), which is
preserved as a potential future supplement for extending \texttt{Dim\_Airline}.

% ===========================================================================
\section{ETL Approach}
% ===========================================================================

\subsection{Chosen Approach}

A \textbf{self-developed Python ETL} (Approach~2 from the laboratory
specification) was selected. The pipeline is implemented as four modules:

\begin{itemize}
    \item \texttt{etl/extract.py} — reads raw data from CSV files.
    \item \texttt{etl/transform.py} — builds dimension DataFrames and
          transforms raw flight records into fact rows.
    \item \texttt{etl/load.py} — writes transformed data to PostgreSQL.
    \item \texttt{etl/etl.py} — the orchestrator that invokes the above
          in the correct order and manages database connections.
\end{itemize}

\subsection{Rationale}

\begin{itemize}
    \item Full control over transformation logic, NULL handling, and
          memory-efficient chunk iteration.
    \item The \texttt{pandas} + \texttt{psycopg2} combination is well-suited
          to the CSV-to-PostgreSQL stack used in this project.
    \item No additional licensed software or services are required.
    \item The code is straightforward to inspect, debug, and extend.
\end{itemize}

An ELT approach (Approach~3) was considered but rejected: it would require
loading approximately 12~GB of raw CSV data into PostgreSQL before any
transformations, which is less efficient than the chunked in-memory Python
approach. A dedicated ETL tool (Approach~1, e.g.\ SSIS) was rejected because
it would introduce a heavyweight dependency for a relatively straightforward
pipeline.

\subsection{Technology Stack}

\begin{center}
\begin{tabular}{lll}
\toprule
\textbf{Component} & \textbf{Tool} & \textbf{Version} \\
\midrule
Extraction     & \texttt{pandas.read\_csv} with chunk iteration & pandas $\geq$ 2.0 \\
Transformation & Python, pandas in-memory DataFrames            & Python 3.9+ \\
Loading        & \texttt{psycopg2} \texttt{execute\_values} + \texttt{COPY FROM STDIN} & psycopg2-binary $\geq$ 2.9 \\
Database       & PostgreSQL 14 running in a Docker container    & PostgreSQL 14 \\
Orchestration  & \texttt{etl/etl.py}                            & — \\
\bottomrule
\end{tabular}
\end{center}

\subsection{Pipeline Overview}

The \texttt{run\_etl()} function in \texttt{etl.py} executes the following
eight steps in order:

\begin{enumerate}
    \item Create (or recreate) the schema from \texttt{sql/ddl.sql}.
    \item Build and load \texttt{Dim\_Date} (programmatically generated).
    \item Extract, transform, and load \texttt{Dim\_Airline}.
    \item Extract, transform, and load \texttt{Dim\_Airport}.
    \item Scan \texttt{flights.csv} for unique tail numbers; build and load
          \texttt{Dim\_Plane}.
    \item Query all dimension surrogate keys into in-memory lookup
          dictionaries.
    \item Stream Dataset~3 flight records in 100\,000-row chunks through the
          transformation pipeline and bulk-load into the fact table.
    \item Pre-scan Dataset~1 for missing airlines; stream flight records
          through the same pipeline and load into the fact table.
\end{enumerate}

% ===========================================================================
\section{Data Staging Area}
% ===========================================================================

No separate staging area database was used. The raw CSV files on disk serve
directly as the source layer. The ETL reads data in memory-efficient chunks
of 100\,000 rows without first materialising the full dataset in any
intermediate database.

This decision is justified because:

\begin{itemize}
    \item The datasets are \emph{static snapshots}, not live feeds — there is
          no ongoing requirement to merge incremental loads.
    \item The transformation logic is simple enough to be applied entirely
          in-memory during chunk iteration.
    \item Loading ~12~GB into a staging database first would double the total
          disk I/O for no analytical benefit.
\end{itemize}

If the pipeline were extended to support incremental loads from live sources,
introducing a staging schema (e.g.\ raw tables mirroring the CSV structure)
would become appropriate.

% ===========================================================================
\section{Transformations}
% ===========================================================================

This section describes every transformation applied prior to loading data into
the data warehouse, based on the issues identified in Section~1.

% ---------------------------------------------------------------------------
\subsection{Dim\_Date — Programmatic Generation}
% ---------------------------------------------------------------------------

\texttt{Dim\_Date} is \emph{not} derived from any source file. It is
generated entirely in Python by the \texttt{build\_dim\_date()} function in
\texttt{transform.py}. For each year in the set
$\{2015, 2019, 2020, 2021, 2022, 2023\}$, every calendar day is enumerated
and the following attributes are computed:

\begin{center}
\begin{tabular}{ll}
\toprule
\textbf{Attribute} & \textbf{Derivation} \\
\midrule
\texttt{date\_key}    & $\text{YEAR} \times 10000 + \text{MONTH} \times 100 + \text{DAY}$ \\
\texttt{full\_date}   & Python \texttt{date} object \\
\texttt{year}         & \texttt{date.year} \\
\texttt{month}        & \texttt{date.month} \\
\texttt{month\_name}  & Static lookup (e.g.\ $1 \to$ \texttt{'January'}) \\
\texttt{quarter}      & $(\text{month} - 1) \div 3 + 1$ \\
\texttt{day\_of\_month} & \texttt{date.day} \\
\texttt{day\_of\_week}  & \texttt{date.isoweekday()} — 1 = Monday, 7 = Sunday \\
\texttt{day\_name}    & Static lookup (e.g.\ $3 \to$ \texttt{'Wednesday'}) \\
\texttt{is\_weekend}  & \texttt{day\_of\_week} $\geq 6$ \\
\texttt{season}       & Month-to-season map: Dec/Jan/Feb = Winter, Mar–May = Spring,\\
                      & Jun–Aug = Summer, Sep–Nov = Autumn \\
\bottomrule
\end{tabular}
\end{center}

This produces 2\,192 rows covering all days across the six years.

% ---------------------------------------------------------------------------
\subsection{Dim\_Airline}
% ---------------------------------------------------------------------------

\textbf{Source:} \texttt{data/dataset3/airlines.csv} (14 rows), supplemented
by carrier codes discovered in Dataset~1 during Step~8 of the pipeline.

\textbf{Transformations applied} (in \texttt{transform\_dim\_airline()}):

\begin{itemize}
    \item Strip leading/trailing whitespace from \texttt{IATA\_CODE} and
          \texttt{AIRLINE}.
    \item Rename columns to the internal schema names
          (\texttt{iata\_code}, \texttt{airline\_name}).
    \item Set \texttt{country\_of\_origin = 'USA'} for all rows (all carriers
          in the dataset are US domestic operators).
    \item Deduplicate by \texttt{iata\_code} to prevent duplicate inserts.
\end{itemize}

\textbf{Dataset~1 supplement:} Before loading Dataset~1 fact rows, the
\texttt{\_ensure\_dataset1\_airlines()} function scans the
\texttt{AIRLINE\_CODE} column in chunks and inserts any IATA codes not yet
present in \texttt{Dim\_Airline}, using the \texttt{AIRLINE} full-name column
as the carrier name. Airline names longer than 50 characters are truncated to
fit the \texttt{VARCHAR(50)} column constraint.

% ---------------------------------------------------------------------------
\subsection{Dim\_Airport}
% ---------------------------------------------------------------------------

\textbf{Source:} \texttt{data/dataset3/airports.csv} (322 rows).

\textbf{Transformations applied} (in \texttt{transform\_dim\_airport()}):

\begin{itemize}
    \item Strip whitespace from all string columns
          (\texttt{IATA\_CODE}, \texttt{AIRPORT}, \texttt{CITY},
          \texttt{COUNTRY}).
    \item Truncate \texttt{airport\_name} to 100 characters (the longest
          observed name is 78 characters, so no data is lost in practice).
    \item Convert \texttt{LATITUDE} and \texttt{LONGITUDE} to
          \texttt{DECIMAL(9,5)} using \texttt{pd.to\_numeric(errors='coerce')};
          any non-numeric value is coerced to \texttt{NULL}.
    \item Deduplicate by \texttt{iata\_code}.
\end{itemize}

% ---------------------------------------------------------------------------
\subsection{Dim\_Plane}
% ---------------------------------------------------------------------------

\textbf{Source:} The \texttt{TAIL\_NUMBER} column scanned from
\texttt{data/dataset3/flights.csv} in a single pass of 100\,000-row chunks.

\textbf{Transformations applied} (in \texttt{build\_dim\_plane()}):

\begin{itemize}
    \item Extract all unique, non-null, non-empty tail number strings
          ($\approx$4\,800 distinct values).
    \item Prepend a special \texttt{UNKNOWN} sentinel row as the first entry.
          This row acts as a foreign-key target for all fact rows that lack a
          tail number.
    \item All metadata columns (\texttt{model}, \texttt{manufacturer},
          \texttt{issue\_date}) are set to \texttt{NULL} because the source
          dataset provides no aircraft metadata beyond the tail number itself.
    \item \texttt{status} is set to the string \texttt{'Unknown'} for all rows.
\end{itemize}

% ---------------------------------------------------------------------------
\subsection{Dim\_Cancellation\_Reason}
% ---------------------------------------------------------------------------

This dimension is seeded \emph{directly in the DDL} with five static values
and requires no ETL step:

\begin{center}
\begin{tabular}{ll}
\toprule
\textbf{Code} & \textbf{Description} \\
\midrule
\texttt{A} & Carrier \\
\texttt{B} & Weather \\
\texttt{C} & National Air System \\
\texttt{D} & Security \\
\texttt{N} & Not Applicable (sentinel for non-cancelled flights) \\
\bottomrule
\end{tabular}
\end{center}

% ---------------------------------------------------------------------------
\subsection{Column Normalisation (Unified Schema)}
% ---------------------------------------------------------------------------

Because Datasets~1 and~3 use different column names for the same concepts,
each chunk is first passed through a normalisation function before entering
the shared fact transformation logic.

\texttt{normalize\_dataset3\_chunk()} (in \texttt{transform.py}) renames the
Dataset~3 columns:

\begin{center}
\begin{tabular}{ll}
\toprule
\textbf{Source column} & \textbf{Internal name} \\
\midrule
\texttt{AIRLINE}             & \texttt{airline\_iata} \\
\texttt{ORIGIN\_AIRPORT}     & \texttt{origin\_iata} \\
\texttt{DESTINATION\_AIRPORT}& \texttt{dest\_iata} \\
\texttt{DEPARTURE\_DELAY}    & \texttt{dep\_delay} \\
\texttt{ARRIVAL\_DELAY}      & \texttt{arr\_delay} \\
\texttt{AIR\_SYSTEM\_DELAY}  & \texttt{nas\_delay} \\
\texttt{AIRLINE\_DELAY}      & \texttt{carrier\_delay} \\
\texttt{CANCELLATION\_REASON}& \texttt{cancel\_code} \\
\ldots & \ldots \\
\bottomrule
\end{tabular}
\end{center}

\texttt{normalize\_dataset1\_chunk()} applies an analogous renaming for
Dataset~1 (e.g.\ \texttt{AIRLINE\_CODE} $\to$ \texttt{airline\_iata},
\texttt{DELAY\_DUE\_CARRIER} $\to$ \texttt{carrier\_delay}) and additionally
sets \texttt{tail\_number = None} for all rows (Dataset~1 has no tail number
column) and copies \texttt{FL\_DATE} into an auxiliary
\texttt{flight\_date\_str} column for date parsing.

% ---------------------------------------------------------------------------
\subsection{Fact Table Transformations (Per-Chunk)}
% ---------------------------------------------------------------------------

Each 100\,000-row chunk passes through \texttt{transform\_fact\_chunk()} in
\texttt{transform.py}. The steps are as follows:

\subsubsection*{1. Date Key Construction}

For Dataset~3, the integer date key is derived directly from the source
columns:
\[
  \texttt{date\_key} = \texttt{YEAR} \times 10000 + \texttt{MONTH} \times 100 + \texttt{DAY}
\]

For Dataset~1, the date string \texttt{FL\_DATE} (\texttt{YYYY-MM-DD}) is
split on the hyphen separator and the same formula is applied. Rows whose
date key does not correspond to a date already loaded in \texttt{Dim\_Date}
are silently dropped.

\subsubsection*{2. Dimension Key Resolution}

Surrogate keys for all four foreign-key dimensions are resolved using
in-memory Python dictionaries (loaded in Step~6 of the pipeline). IATA codes
are stripped of whitespace before lookup:

\begin{itemize}
    \item \textbf{Airline key:} \texttt{airline\_iata} $\to$ \texttt{airline\_key}. Rows
          with unresolved airline codes are dropped.
    \item \textbf{Airport keys:} \texttt{origin\_iata} and \texttt{dest\_iata}
          $\to$ \texttt{origin\_airport\_key} and
          \texttt{destination\_airport\_key}. Rows where either airport
          cannot be resolved are counted as skipped and excluded from the
          output.
    \item \textbf{Plane key:} \texttt{tail\_number} $\to$ \texttt{plane\_key}.
          Any tail number not in the lookup (including explicit \texttt{None},
          empty string, or the string \texttt{'nan'}) is mapped to the
          \texttt{UNKNOWN} sentinel key.
    \item \textbf{Cancellation reason key:} resolved from the single-character
          \texttt{cancel\_code}. For non-cancelled flights (code absent or not
          in the lookup), the \texttt{N} sentinel key is used.
\end{itemize}

\subsubsection*{3. HHMM Time Parsing}

The helper function \texttt{parse\_hhmm()} converts any HHMM representation
(integer, float, or string) to the format \texttt{HH:MM:00} expected by
PostgreSQL's \texttt{TIME} type:

\begin{itemize}
    \item Float decimals are stripped (e.g.\ \texttt{930.0} $\to$ \texttt{'0930'}).
    \item The string is zero-padded to four characters
          (\texttt{'930'} $\to$ \texttt{'0930'}).
    \item The edge case \texttt{'2400'} (midnight expressed as end-of-day) is
          mapped to \texttt{'00:00:00'}.
    \item Values with out-of-range hour ($>23$) or minute ($>59$) are
          returned as \texttt{None}.
    \item Missing or \texttt{NaN} values return \texttt{None}.
\end{itemize}

\subsubsection*{4. NULL and Zero Handling for Delay Columns}

\begin{itemize}
    \item Cancelled flights: \texttt{dep\_delay} and \texttt{arr\_delay} are
          stored as \texttt{NULL} (flight never departed or arrived).
    \item Non-delayed, non-cancelled flights: all five delay breakdown columns
          (\texttt{carrier\_delay}, \texttt{weather\_delay},
          \texttt{nas\_delay}, \texttt{security\_delay},
          \texttt{late\_aircraft\_delay}) are stored as \texttt{0}.
    \item Delayed flights: breakdown columns retain their original values.
\end{itemize}

\subsubsection*{5. Type Casting}

All delay, time, and distance columns are cast from pandas
\texttt{Float32}/\texttt{Float64} to Python \texttt{int} (or \texttt{None}
for \texttt{NaN}) via the \texttt{\_to\_int\_or\_none()} helper. This is
required because PostgreSQL's \texttt{SMALLINT} columns do not accept
floating-point values.

\subsubsection*{6. Source Label}

A \texttt{source\_dataset} string (\texttt{'dataset3'} or
\texttt{'dataset1'}) is added to every row. This column is not present in
either source file; it is injected by the orchestrator to allow downstream
queries to filter or aggregate by data vintage.

% ===========================================================================
\section{Loading Phase}
% ===========================================================================

\subsection{Schema Creation}

At the start of every ETL run, \texttt{create\_schema()} reads
\texttt{sql/ddl.sql} and executes it against the target database. All tables
are dropped in reverse foreign-key order (\texttt{CASCADE}) and recreated,
ensuring a clean state. This makes the full ETL run idempotent and
reproducible.

\subsection{Dimension Loading}

All dimension tables are loaded using
\texttt{psycopg2.extras.execute\_values()}, which batches multiple row tuples
into a single parameterised \texttt{INSERT} statement per call, reducing
network round-trips compared to per-row execution.

Every dimension insert uses the clause \texttt{ON CONFLICT DO NOTHING},
meaning:

\begin{itemize}
    \item The ETL can be safely re-run without duplicating dimension data.
    \item Partial runs (e.g.\ interrupted after dimension loading) do not
          cause constraint violations on retry.
\end{itemize}

Loading order respects the absence of cross-dimension foreign-key
dependencies:

\begin{enumerate}
    \item \texttt{Dim\_Date} — 2\,192 rows (programmatically generated)
    \item \texttt{Dim\_Airline} — 14+ rows (from \texttt{airlines.csv} plus Dataset~1 supplements)
    \item \texttt{Dim\_Airport} — 322 rows (from \texttt{airports.csv})
    \item \texttt{Dim\_Plane} — $\approx$4\,800 rows (scanned from \texttt{flights.csv})
    \item \texttt{Dim\_Cancellation\_Reason} — 5 rows (pre-seeded by DDL; no ETL step)
\end{enumerate}

\subsection{Lookup Dictionary Build}

After all dimension tables are populated, \texttt{build\_dimension\_lookup\_dicts()} issues five \texttt{SELECT} queries and constructs
in-memory Python dictionaries mapping each business key (IATA code, tail
number, cancellation code, date key) to its surrogate key. These dictionaries
are used throughout fact loading for O(1) key resolution without issuing a
database query per row.

\subsection{Fact Table Loading — COPY FROM STDIN}

The fact table is loaded using PostgreSQL's native
\texttt{COPY FROM STDIN} protocol, accessed through
\texttt{psycopg2}'s \texttt{cursor.copy\_expert()} method. The process
per chunk is:

\begin{enumerate}
    \item A transformed list of tuples is serialised to a tab-separated
          in-memory \texttt{io.StringIO} buffer by the \texttt{\_row\_to\_tsv()}
          helper. Python \texttt{None} values are emitted as \texttt{{\textbackslash}N}
          (the PostgreSQL text-format null marker). Boolean values are written
          as \texttt{t} or \texttt{f}.
    \item The buffer is rewound and passed to \texttt{copy\_expert()} with the
          following \texttt{COPY} statement:
\begin{lstlisting}[style=sql]
COPY Fact_Flight_Operations (
    date_key, airline_key, origin_airport_key, destination_airport_key,
    plane_key, cancel_reason_key, flight_number, source_dataset,
    departure_delay_min, arrival_delay_min,
    carrier_delay_min, weather_delay_min, nas_delay_min,
    security_delay_min, late_aircraft_delay_min,
    taxi_out_min, taxi_in_min, air_time_min, distance_miles,
    flight_count, cancelled_flag, diverted_flag
) FROM STDIN WITH (FORMAT TEXT, NULL '\N')
\end{lstlisting}
    \item A \texttt{COMMIT} is issued after each chunk, so progress is
          preserved even if the process is interrupted mid-run.
\end{enumerate}

The performance advantage of \texttt{COPY} over batched \texttt{INSERT} is
substantial:

\begin{center}
\begin{tabular}{lc}
\toprule
\textbf{Loading method} & \textbf{Estimated time for $\sim$5.8M rows} \\
\midrule
\texttt{execute\_values} (batched INSERT) & 2–4 hours \\
\texttt{COPY FROM STDIN}                  & 10–20 minutes \\
\bottomrule
\end{tabular}
\end{center}

\subsection{Split and Join Operations Required by the Warehouse Schema}

The source data schemas differ significantly from the star schema in the data
warehouse. The following split/join operations were performed during loading:

\begin{center}
\begin{tabular}{>{\raggedright\arraybackslash}p{4.5cm}
                >{\raggedright\arraybackslash}p{9.5cm}}
\toprule
\textbf{Operation} & \textbf{Description} \\
\midrule
Role-playing airport dimension &
A single \texttt{Dim\_Airport} table is referenced twice from the fact table
— once as \texttt{origin\_airport\_key} and once as
\texttt{destination\_airport\_key}. Both keys are resolved from the same
in-memory dictionary using the origin and destination IATA codes
respectively. \\[6pt]

Date reconstruction from split columns &
Dataset~3 stores the date across three separate columns
(\texttt{YEAR}, \texttt{MONTH}, \texttt{DAY}).
These are combined into a single integer \texttt{date\_key} during
transformation. \\[6pt]

Airline surrogate key join &
Source files contain only the IATA carrier code string. The integer
\texttt{airline\_key} surrogate is resolved via the in-memory dictionary. \\[6pt]

Plane surrogate key join &
Tail number strings are mapped to integer \texttt{plane\_key} values. Rows
with missing or unrecognised tail numbers are mapped to the \texttt{UNKNOWN}
sentinel key. \\[6pt]

Cancellation reason join &
The raw single-character cancellation code is mapped to the integer
\texttt{cancel\_reason\_key}; non-cancelled rows receive the \texttt{N}
sentinel. \\[6pt]

Source dataset label injection &
The \texttt{source\_dataset} column (\texttt{'dataset3'} or
\texttt{'dataset1'}) does not exist in any source file and is injected by the
ETL orchestrator, enabling downstream queries to filter by data vintage or
combine both sources with explicit awareness of the difference. \\

\bottomrule
\end{tabular}
\end{center}

\subsection{Final Row Counts}

After a complete ETL run, the warehouse contains approximately:

\begin{center}
\begin{tabular}{lrl}
\toprule
\textbf{Table} & \textbf{Rows} & \textbf{Source} \\
\midrule
\texttt{Dim\_Date}                & 2\,192       & Programmatically generated \\
\texttt{Dim\_Airline}             & 14+          & \texttt{airlines.csv} + Dataset~1 supplement \\
\texttt{Dim\_Airport}             & 322          & \texttt{airports.csv} \\
\texttt{Dim\_Plane}               & $\sim$4\,800 & Scanned from \texttt{flights.csv} \\
\texttt{Dim\_Cancellation\_Reason}& 5            & DDL seed \\
\texttt{Fact\_Flight\_Operations} & $\sim$8\,800\,000 & Dataset~3 ($\sim$5.8M) + Dataset~1 ($\sim$3M) \\
\bottomrule
\end{tabular}
\end{center}

\end{document}
